\chapter{2011年辽宁卷（文）}

\section{单选题}

\begin{problem}
已知集合 \(A = \{x \mid x > 1\}\)，\(B = \{x \mid -1 < x < 2\}\)，则 \(A \cap B =\)（\quad）

\choices
    {\(\{x\mid -1 < x < 2\}\)}
    {\( \{x \mid x > -1\} \)}
    {\( \{x \mid -1 < x < 1\} \)}
    {\( \{x \mid 1 < x < 2\} \)}
\end{problem}

\begin{answer}
D
\end{answer}

\begin{solution}
由交集定义，元素必须同时满足 \(x>1\) 与 \(-1<x<2\)，所以 \(1<x<2\).因而 \(A\cap B=\{x\mid 1<x<2\}\)，故选 D.
\end{solution}

\begin{problem}
\(\i\) 为虚数单位，\(\frac{1}{\i} + \frac{1}{\i^{3}} + \frac{1}{\i^{5}} + \frac{1}{\i^{7}} =\)（\quad）

\choices
    {\(0\)}
    {\(2\i\)}
    {\(-2\i\)}
    {\(4\i\)}
\end{problem}

\begin{answer}
A
\end{answer}

\begin{solution}
因为 \(\i^2=-1\)，所以 \(\i^3=-\i\)，\(\ \i^5=\i\)，\(\ \i^7=-\i\).于是
\(\displaystyle \frac{1}{\i}+\frac{1}{\i^3}+\frac{1}{\i^5}+\frac{1}{\i^7}
=-\i+\i-\i+\i=0\)，故选 A.
\end{solution}

\begin{problem}
已知向量 \(\bs{a} = (2,1)\)，\(\bs{b} = (-1,k)\)，\(\bs{a} \cdot (2\bs{a} - \bs{b}) = 0\)，则 \(k =\)（\quad）

\choices
    {\(-12\)}
    {\(-6\)}
    {\(6\)}
    {\(12\)}
\end{problem}

\begin{answer}
D
\end{answer}

\begin{solution}
由 \(2\bs a-\bs b=(4,2)-(-1,k)=(5,2-k)\)，得
\(\bs a\cdot(2\bs a-\bs b)=(2,1)\cdot(5,2-k)=10+2-k=12-k\).题设内积为 \(0\)，故 \(k=12\)，选 D.
\end{solution}

\begin{problem}
已知命题 \(p: \exists n \in \N\)，\(2^n > 1000\)，则 \(\neg p\) 为（\quad）

\choices
    {\(\forall n\in \N,2^{n}\leqslant 1000\)}
    {\(\forall n\in \N\)，\(2^n >1000\)}
    {\(\exists n\in \N\)，\(2^n\leqslant 1000\)}
    {\(\exists n\in \N\)，\(2^n < 1000\)}
\end{problem}

\begin{answer}
A
\end{answer}

\begin{solution}
命题“存在 \(n\in\N\)，使 \(2^n>1000\)”的否定是“对任意 \(n\in\N\)，都有 \(2^n\leqslant1000\)”.故选 A.
\end{solution}

\begin{problem}
若等比数列 \(\{a_{n}\}\) 满足 \(a_{n}a_{n + 1} = 16^{n}\)，则公比为（\quad）

\choices
    {\(2\)}
    {\(4\)}
    {\(8\)}
    {\(16\)}
\end{problem}

\begin{answer}
B
\end{answer}

\begin{solution}
设等比数列的公比为 \(q\).由 \(a_n=a_1q^{n-1}\)，有
\(\displaystyle a_na_{n+1}=a_1^2q^{2n-1}=16^n\).分别写出 \(n+1\) 和 \(n\) 时的等式并相除，得 \(q^2=16\).再由 \(n=1\) 时 \(a_1^2q=16>0\)，且 \(a_1\neq0\)，可知 \(q>0\)，所以 \(q=4\)，选 B.
\end{solution}

\begin{problem}
若函数 \( f(x) = \frac{x}{(2x + 1)(x - a)} \) 为奇函数，则 \( a = \)（\quad）

\choices
    {\(\frac{1}{2}\)}
    {\( \frac{2}{3} \)}
    {\(\frac{3}{4}\)}
    {\(1\)}
\end{problem}

\begin{answer}
A
\end{answer}

\begin{solution}
分母的零点为 \(-\tfrac12\) 和 \(a\).奇函数的定义域关于原点对称，因此
\(\{-\tfrac12,a\}=\{\tfrac12,-a\}\)，从而 \(a=\tfrac12\).此时
\(\displaystyle (2x+1)(x-\tfrac12)=2x^2-\tfrac12\) 是偶函数，故
\(\displaystyle f(-x)=\frac{-x}{2x^2-\tfrac12}=-f(x)\).所以 \(f(x)\) 为奇函数，选 A.
\end{solution}

\begin{problem}
已知 \(F\) 是抛物线 \(y^{2}=x\) 的焦点，\(A\)，\(B\) 是该抛物线上的两点，\(|AF|+|BF|=3\)，则线段 \(AB\) 的中点到 \(y\) 轴的距离为（\quad）

\choices
    {\( \frac{3}{4} \)}
    {\(1\)}
    {\( \frac{5}{4} \)}
    {\( \frac{7}{4} \)}
\end{problem}

\begin{answer}
C
\end{answer}

\begin{solution}
抛物线 \(y^2=x\) 可写为 \(y^2=4px\)，所以 \(p=\tfrac14\)，焦点为 \(F(\tfrac14,0)\).设 \(A(x_1,y_1)\)、\(B(x_2,y_2)\) 在抛物线上，则
\(\displaystyle |AF|^2=(x_1-\tfrac14)^2+y_1^2=(x_1+\tfrac14)^2\)，且 \(x_1\geq0\)，从而 \(|AF|=x_1+\tfrac14\)，同理 \(|BF|=x_2+\tfrac14\).由 \(|AF|+|BF|=3\)，得 \(x_1+x_2=\tfrac52\).弦 \(AB\) 中点的横坐标为 \(\tfrac54\)，所以其中点到 \(y\) 轴的距离为 \(\tfrac54\)，故选 C.
\end{solution}

\begin{problem}
一个正三棱柱的侧棱长和底面边长相等，体积为 \(2\sqrt{3}\)，它的三视图中的俯视图如图所示，左视图是一个矩形，则这个矩形的面积是（\quad）

\bitmapfigure[max width=0.36\linewidth]{img/2011/liaoning_liberal/q8_fig1.png}

\choices
    {\(4\)}
    {\(2\sqrt{3}\)}
    {\(2\)}
    {\(\sqrt{3}\)}
\end{problem}

\begin{answer}
B
\end{answer}

\begin{solution}
设正三棱柱的底面边长为 \(a\)，则侧棱长和高也为 \(a\).底面是边长为 \(a\) 的正三角形，面积为 \(\tfrac{\sqrt3}{4}a^2\)，所以
\(\displaystyle 2\sqrt3=V=\frac{\sqrt3}{4}a^3\)，得 \(a=2\).俯视图给出底面正三角形，左视图的两条边分别为高 \(2\) 和正三角形的高 \(\sqrt3\)，故矩形面积为 \(2\sqrt3\)，选 B.
\end{solution}

\begin{problem}
执行下面的程序框图，如果输入的 \(n\) 是 4，则输出的 \(p\) 是（\quad）

\bitmapfigure[max width=0.52\linewidth]{img/2011/liaoning_liberal/q10_fig1.png}

\choices
    {\(8\)}
    {\(5\)}
    {\(3\)}
    {\(2\)}
\end{problem}

\begin{answer}
C
\end{answer}

\begin{solution}
初始 \(s=0\)，\(t=1\)，\(k=1\).第一次循环 \(p=s+t=1\)，更新为 \((s,t,k)=(1,1,2)\)；第二次循环 \(p=2\)，更新为 \((1,2,3)\)；第三次循环 \(p=3\)，更新为 \((2,3,4)\).此时 \(k<n\) 不成立，输出 \(p=3\)，故选 C.
\end{solution}

\begin{problem}
已知球的直径 \(SC = 4\)，\(A\)，\(B\) 是该球球面上的两点，\(AB = 2\)，\(\angle ASC = \angle BSC = 45^\circ\)，则棱锥 \(S-ABC\) 的体积为（\quad）

\choices
    {\(\frac{\sqrt{3}}{3}\)}
    {\(\frac{2\sqrt{3}}{3}\)}
    {\(\frac{4\sqrt{3}}{3}\)}
    {\(\frac{5\sqrt{3}}{3}\)}
\end{problem}

\begin{answer}
C
\end{answer}

\begin{solution}
取 \(S=(0,0,0)\)、\(C=(4,0,0)\).因为 \(SC\) 是球的直径，\(\angle SAC=90^\circ\)；又 \(\angle ASC=45^\circ\)，所以 \(\triangle SAC\) 为等腰直角三角形，\(SA=AC=2\sqrt2\)，且 \(A\) 在 \(SC\) 上的投影横坐标为 \(2\)，到 \(SC\) 的距离为 \(2\).点 \(B\) 同样满足这两个条件.在垂直于 \(SC\) 的平面内建立直角坐标，使 \(A=(2,2,0)\)，设 \(B=(2,u,v)\).由 \(B\) 到 \(SC\) 的距离为 \(2\) 及 \(AB=2\)，有
\(\displaystyle u^2+v^2=4\)，\((u-2)^2+v^2=4\)，两式相减得 \(u=1\)，再得 \(v=\pm\sqrt3\).旋转坐标方向可取 \(v=\sqrt3\)，故 \(B=(2,1,\sqrt3)\).因此
\(\displaystyle V_{S-ABC}=\frac16\left|\det\begin{pmatrix}2&2&0\\2&1&\sqrt3\\4&0&0\end{pmatrix}\right|
=\frac{4\sqrt3}{3}\)，故选 C.
\end{solution}

\begin{problem}
\(f(x)\) 的定义域为 \(\R\)，\(f(-1) = 2\)，对任意 \(x \in \R\)，\(f'(x) > 2\)，则 \(f(x) > 2x + 4\) 的解集为（\quad）

\choices
    {\((-1,1)\)}
    {\((-1,+\infty)\)}
    {\((-\infty,-1)\)}
    {\((-\infty,+\infty)\)}
\end{problem}

\begin{answer}
B
\end{answer}

\begin{solution}
令 \(g(x)=f(x)-(2x+4)\)，则 \(g'(x)=f'(x)-2>0\)，所以 \(g\) 在 \(\R\) 上严格递增.又
\(\displaystyle g(-1)=f(-1)-2(-1)-4=2+2-4=0\).因此 \(g(x)>0\) 当且仅当 \(x>-1\)，即 \(f(x)>2x+4\) 的解集为 \((-1,+\infty)\)，选 B.
\end{solution}

\begin{problem}
已知函数 \(f(x) = A \tan(\omega x + \varphi)\)（\(\omega > 0\)，\(|\varphi| < \frac{\pi}{2}\)），\(y = f(x)\) 的部分图象如图，则 \(f\left(\frac{\pi}{24}\right) =\)（\quad）

\bitmapfigure[max width=0.34\linewidth]{img/2011/liaoning_liberal/q12_fig1.png}

\choices
    {\(2 + \sqrt{3}\)}
    {\(\sqrt{3}\)}
    {\(\frac{\sqrt{3}}{3}\)}
    {\(2 - \sqrt{3}\)}
\end{problem}

\begin{answer}
B
\end{answer}

\begin{solution}
由图可知，\(x=\frac{\pi}{8}\) 是一条竖直渐近线，\(x=\frac{3\pi}{8}\) 是图像与 \(x\) 轴的交点.正切函数从渐近线到零点的水平距离为 \(\tfrac{\pi}{2\omega}\)，因而
\(\displaystyle \frac{3\pi}{8}-\frac{\pi}{8}=\frac{\pi}{4}=\frac{\pi}{2\omega}\)，得 \(\omega=2\).又在 \(x=\frac{\pi}{8}\) 处有 \(2\cdot\frac{\pi}{8}+\varphi=\frac{\pi}{2}\) 的渐近线，结合 \(|\varphi|<\frac{\pi}{2}\) 得 \(\varphi=\frac{\pi}{4}\).由图知 \(f(0)=1\)，故 \(A\tan\frac{\pi}{4}=A=1\).因此
\(\displaystyle f\left(\frac{\pi}{24}\right)=\tan\left(2\cdot\frac{\pi}{24}+\frac{\pi}{4}\right)
=\tan\frac{\pi}{3}=\sqrt3\)，故选 B.
\end{solution}

\section{填空题}

\begin{problem}
已知圆 \(C\) 经过 \(A(5,1)\)，\(B(1,3)\) 两点，圆心在 \(x\) 轴上，则 \(C\) 的方程为\fillinblank{}.
\end{problem}

\begin{answer}
\((x-2)^2+y^2=10\)
\end{answer}

\begin{solution}
设圆心为 \((h,0)\)，半径为 \(r\).由 \(A(5,1)\)、\(B(1,3)\) 在圆上，
\(\displaystyle (5-h)^2+1=(1-h)^2+9\).展开化简得 \(h=2\)，再代入得
\(\displaystyle r^2=(5-2)^2+1=10\).所以圆的方程为 \((x-2)^2+y^2=10\).
\end{solution}

\begin{problem}
调查了某地若干户家庭的年收入 \(x\)（单位：万元）和年饮食支出 \(y\)（单位：万元）. 调查显示年收入 \(x\) 与年饮食支出 \(y\) 具有线性相关关系，并由调查数据得到 \(y\) 对 \(x\) 的回归直线方程：\(\hat{y} = 0.254x + 0.321\). 由回归直线方程可知，家庭年收入每增加 \(1\) 万元，年饮食支出平均增加\fillinblank{} 万元.
\end{problem}

\begin{answer}
\(0.254\)
\end{answer}

\begin{solution}
回归直线 \(\hat y=0.254x+0.321\) 的斜率为 \(0.254\).收入 \(x\) 增加 \(1\) 万元时，预测饮食支出的增加量为斜率，即 \(0.254\) 万元.
\end{solution}

\begin{problem}
\(S_{n}\) 为等差数列 \(\{a_{n}\}\) 的前 \(n\) 项和，\(S_{2} = S_{6}\)，\(a_{4} = 1\)，则 \(a_{5} =\)\fillinblank{}.
\end{problem}

\begin{answer}
\(-1\)
\end{answer}

\begin{solution}
设等差数列首项为 \(a_1\)，公差为 \(d\).由
\(\displaystyle S_2=2a_1+d\)，
\(\displaystyle S_6=\frac62(2a_1+5d)=6a_1+15d\)，以及 \(S_2=S_6\)，得
\(\displaystyle 2a_1+7d=0\).又 \(a_4=a_1+3d=1\).代入 \(a_1=-\tfrac72d\)，得 \(-\tfrac12d=1\)，所以 \(d=-2\)，\(a_1=7\).因此 \(a_5=a_1+4d=7-8=-1\).
\end{solution}

\begin{problem}
已知函数 \(f(x) = \e^{x} - 2x + a\) 有零点，则 \(a\) 的取值范围是\fillinblank{}.
\end{problem}

\begin{answer}
\((-\infty,2\ln 2-2]\)
\end{answer}

\begin{solution}
令 \(g(x)=\e^x-2x\)，则 \(g'(x)=\e^x-2\).函数在 \(x=\ln2\) 左侧递减、右侧递增，故最小值为
\(\displaystyle g(\ln2)=2-2\ln2\).又 \(g(x)\to+\infty\) 当 \(x\to\pm\infty\)，所以 \(\e^x-2x+a=0\) 有实根当且仅当最小值不大于 \(0\)，即
\(\displaystyle 2-2\ln2+a\leq0\).故 \(a\in(-\infty,2\ln2-2]\).
\end{solution}

\section{解答题}

\begin{problem}
\(\triangle ABC\) 的三个内角 \(A\)，\(B\)，\(C\) 所对的边分别为 \(a\)，\(b\)，\(c\)，\(a \sin A \sin B + b \cos^{2} A = \sqrt{2}a\).
\begin{enumerate}
\item 求 \(\frac{b}{a}\)；
\item 若 \(c^{2} = b^{2} + \sqrt{3}a^{2}\)，求 \(B\).
\end{enumerate}
\end{problem}

\begin{solution}
\begin{enumerate}
\item 由正弦定理，\(\displaystyle \frac{b}{a}=\frac{\sin B}{\sin A}\).将题设等式两边同除以 \(a\)，并代入上式，得
\(\displaystyle \sin A\sin B+\frac{\sin B}{\sin A}\cos^2A
=\frac{\sin B}{\sin A}(\sin^2A+\cos^2A)
=\frac{\sin B}{\sin A}=\sqrt2\).所以 \(\displaystyle \frac ba=\sqrt2\).
\item 由上一问 \(b^2=2a^2\)，故题设给出
\(\displaystyle c^2=(2+\sqrt3)a^2\).由余弦定理，
\(\displaystyle \cos B=\frac{a^2+c^2-b^2}{2ac}
=\frac{(1+\sqrt3)a^2}{2a^2\sqrt{2+\sqrt3}}
=\frac{\sqrt2}{2}\)，其中用了 \((1+\sqrt3)^2=2(2+\sqrt3)\).因 \(B\) 是三角形内角且余弦为正，故 \(B=45^\circ\).
\end{enumerate}
\end{solution}

\begin{problem}
如图，四边形 \(ABCD\) 为正方形，\(QA \perp\) 平面 \(ABCD\)，\(PD \parallel QA\)，\(QA = AB = \frac{1}{2}PD\).
\begin{enumerate}
\item 证明：\(PQ \perp\) 平面 \(DCQ\)；
\item 求棱锥 \(Q-ABCD\) 的体积与棱锥 \(P-DCQ\) 的体积比值.
\end{enumerate}

\bitmapfigure[max width=0.38\linewidth]{img/2011/liaoning_liberal/q18_fig1.png}
\end{problem}

\begin{solution}
\begin{enumerate}
\item 令 \(a=AB=QA\)，建立空间直角坐标系，取
\(\displaystyle A=(0,0,0)\)，\(B=(a,0,0)\)，\(C=(a,a,0)\)，\(D=(0,a,0)\)，\(Q=(0,0,a)\).由 \(PD\parallel QA\)、\(PD=2a\) 及图示位置，可取 \(P=(0,a,2a)\).于是
\(\displaystyle \overrightarrow{PQ}=(0,-a,-a)\)，\(\overrightarrow{DC}=(a,0,0)\)，\(\overrightarrow{DQ}=(0,-a,a)\).
从而 \(\overrightarrow{PQ}\cdot\overrightarrow{DC}=0\)，\(\overrightarrow{PQ}\cdot\overrightarrow{DQ}=0\).直线 \(DC\)，\(DQ\) 相交于 \(D\) 且在平面 \(DCQ\) 内，所以 \(PQ\perp\) 平面 \(DCQ\).
\item \(\displaystyle V_{Q-ABCD}=\frac13a^2\cdot a=\frac{a^3}{3}\).又 \(DC=a\)，\(DQ=\sqrt2a\)，所以
\(\displaystyle S_{\triangle DCQ}=\frac12a\cdot\sqrt2a=\frac{\sqrt2}{2}a^2\)，且 \(PQ=\sqrt2a\).因此
\(\displaystyle V_{P-DCQ}=\frac13\cdot\frac{\sqrt2}{2}a^2\cdot\sqrt2a=\frac{a^3}{3}\)，所求体积比为 \(1\).
\end{enumerate}
\end{solution}

\begin{problem}
某农场计划种植某种新作物，为此对这种作物的两个品种（分别称为品种甲和品种乙）进行田间试验，选取两大块地. 每大块地分成 \(n\) 小块地. 在总共 \(2n\) 小块地中，随机选 \(n\) 小块地种植品种甲，另外 \(n\) 小块地种植品种乙.
\begin{enumerate}
\item 假设 \(n = 2\)，求第一大块地都种植品种甲的概率；
\item 试验时每大块地分成 \(8\) 小块，即 \(n = 8\)，试验结束后得到品种甲和品种乙在各小块地上的每公顷产量（单位：\(\mathrm{kg/hm^2}\)）如下表：
\begin{center}
\begin{tabular}{c|cccccccc}
\hline
品种甲 & \(403\) & \(397\) & \(390\) & \(404\) & \(388\) & \(400\) & \(412\) & \(406\) \\
\hline
品种乙 & \(419\) & \(403\) & \(412\) & \(418\) & \(408\) & \(423\) & \(400\) & \(413\) \\
\hline
\end{tabular}
\end{center}
分别求品种甲和品种乙每公顷产量的样本平均数和样本方差；根据试验结果，你应该种植哪一品种？
\end{enumerate}

附：样本数据 \(x_{1}\)，\(x_{2}\)，\(\cdots\)，\(x_{n}\) 的样本方差 \(s^{2} = \frac{1}{n}[(x_{1} - \overline{x})^{2} + (x_{2} - \overline{x})^{2} + \cdots + (x_{n} - \overline{x})^{2}]\)，其中 \(\overline{x}\) 为样本平均数.
\end{problem}

\begin{solution}
\begin{enumerate}
\item 当 \(n=2\) 时，设第一大块地中的两个小块编号为 \(1,2\)，第二大块地中的两个小块编号为 \(3,4\).从四个小块中任选两个种植品种甲，共有 \(\mathrm C_4^2=6\) 种等可能结果；第一大块地都种植品种甲只有选取 \(\{1,2\}\) 这一种结果，所以所求概率为 \(\displaystyle \frac{1}{\mathrm C_4^2}=\frac16\).
\item 品种甲的样本平均数为
\(\displaystyle \overline{x}=\frac{403+397+390+404+388+400+412+406}{8}=400\).
其离均差平方和为
\(\displaystyle 3^2+(-3)^2+(-10)^2+4^2+(-12)^2+0^2+12^2+6^2=458\)，所以样本方差
\(\displaystyle s_x^2=\frac{458}{8}=57.25\).

品种乙的样本平均数为
\(\displaystyle \overline{y}=\frac{419+403+412+418+408+423+400+413}{8}=412\).
其离均差平方和为
\(\displaystyle 7^2+(-9)^2+0^2+6^2+(-4)^2+11^2+(-12)^2+1^2=448\)，所以样本方差
\(\displaystyle s_y^2=\frac{448}{8}=56\).
品种乙的样本平均数 \(412>400\)，且 \(s_y^2<s_x^2\)，说明其平均产量更高、波动更小，故应选择品种乙.
\end{enumerate}
\end{solution}

\begin{problem}
设函数 \(f(x) = x + ax^{2} + b \ln x\)，曲线 \(y = f(x)\) 过 \(P(1,0)\)，且在 \(P\) 点处的切线斜率为 \(2\).
\begin{enumerate}
\item 求 \(a\)，\(b\) 的值；
\item 证明：\(f(x) \leqslant 2x - 2\).
\end{enumerate}
\end{problem}

\begin{solution}
\begin{enumerate}
\item 由 \(f(1)=0\)，得 \(1+a=0\)，所以 \(a=-1\).又
\(\displaystyle f'(x)=1+2ax+\frac{b}{x}\)，由 \(f'(1)=2\) 得 \(1+2a+b=2\)，所以 \(b=3\).
\item 此时 \(f(x)=x-x^2+3\ln x\)，定义域为 \(x>0\).令
\(\displaystyle h(x)=2x-2-f(x)=x^2+x-2-3\ln x\)，则
\(\displaystyle h'(x)=2x+1-\frac3x=\frac{(2x+3)(x-1)}x\).
当 \(0<x<1\) 时 \(h'(x)<0\)，当 \(x>1\) 时 \(h'(x)>0\)，所以 \(h\) 在 \(x=1\) 处取得最小值.又 \(h(1)=0\)，故 \(h(x)\geq0\)，即 \(f(x)\leq2x-2\).
\end{enumerate}
\end{solution}

\begin{problem}
如图，已知椭圆 \(C_{1}\) 的中心在原点 \(O\)，长轴左，右端点 \(M\)，\(N\) 在 \(x\) 轴上，椭圆 \(C_{2}\) 的短轴为 \(MN\)，且 \(C_{1}\)，\(C_{2}\) 的离心率都为 \(e\). 直线 \(l \perp MN\)，\(l\) 与 \(C_{1}\) 交于两点，与 \(C_{2}\) 交于两点，这四点按纵坐标从大到小依次为 \(A\)，\(B\)，\(C\)，\(D\).
\begin{enumerate}
\item 设 \(e = \frac{1}{2}\)，求 \(|BC|\) 与 \(|AD|\) 的比值；
\item 当 \(e\) 变化时，是否存在直线 \(l\)，使得 \(BO \parallel AN\)，并说明理由.
\end{enumerate}

\bitmapfigure[max width=0.38\linewidth]{img/2011/liaoning_liberal/q21_fig1.png}
\end{problem}

\begin{solution}
\begin{enumerate}
\item 设 \(C_1\) 的长半轴、短半轴分别为 \(a\)，\(b\)，则
\(\displaystyle e^2=1-\frac{b^2}{a^2}\).因为 \(MN\) 是 \(C_2\) 的短轴，且两椭圆离心率相同，可设
\(\displaystyle C_1:\frac{x^2}{a^2}+\frac{y^2}{b^2}=1,\quad
C_2:\frac{x^2}{a^2}+\frac{b^2y^2}{a^4}=1\).
设 \(l:x=t\)，其中 \(|t|<a\)，并令 \(A\)，\(B\) 为上方两个交点，则
\(\displaystyle y_A=\frac ab\sqrt{a^2-t^2},\quad
y_B=\frac ba\sqrt{a^2-t^2}\).由关于 \(x\) 轴的对称性，
\(\displaystyle |AD|=2y_A\)，\(|BC|=2y_B\)，故
\(\displaystyle \frac{|BC|}{|AD|}=\frac{b^2}{a^2}=1-e^2=\frac34\).
\item 设 \(O=(0,0)\)，\(N=(a,0)\).当 \(t=0\) 时，\(BO\) 为竖直线而 \(AN\) 不是竖直线，不满足平行；因此取 \(t\neq0\).由
\(\displaystyle B=\left(t,\frac ba\sqrt{a^2-t^2}\right),\quad
A=\left(t,\frac ab\sqrt{a^2-t^2}\right)\)，条件 \(BO\parallel AN\) 等价于斜率相等：
\(\displaystyle \frac{(b/a)\sqrt{a^2-t^2}}{t}
=\frac{(a/b)\sqrt{a^2-t^2}}{t-a}\).
解得
\(\displaystyle t=-\frac{ab^2}{a^2-b^2}
=-a\frac{1-e^2}{e^2}\).要使直线与椭圆有两个交点，须 \(|t|<a\)，这等价于
\(\displaystyle \frac{1-e^2}{e^2}<1\)，即 \(e>\frac{\sqrt2}{2}\).结合 \(0<e<1\)，当 \(0<e\leqslant\frac{\sqrt2}{2}\) 时不存在，当 \(\frac{\sqrt2}{2}<e<1\) 时存在.
\end{enumerate}
\end{solution}

\section{选考题：解答题}

\begin{problem}
如图，\(A\)，\(B\)，\(C\)，\(D\) 四点在同一圆上，\(AD\) 的延长线与 \(BC\) 的延长线交于 \(E\) 点，且 \(EC = ED\).
\begin{enumerate}
\item 证明： \(CD \parallel AB\)；
\item 延长 \(CD\) 到 \(F\)，延长 \(DC\) 到 \(G\)，使得 \(EF = EG\)，证明： \(A\)，\(B\)，\(G\)，\(F\) 四点共圆.
\end{enumerate}

\bitmapfigure[max width=0.34\linewidth]{img/2011/liaoning_liberal/q22_fig1.png}
\end{problem}

\begin{solution}
\begin{enumerate}
\item 因为 \(EC=ED\)，所以 \(\angle EDC=\angle ECD\).又 \(A\)，\(D\)，\(E\) 共线、\(B\)，\(C\)，\(E\) 共线，在圆 \(ABCD\) 中，\(\angle ADC+\angle ABC=180^\circ\)，而 \(\angle EDC\) 与 \(\angle ADC\) 互为补角，\(\angle EBA=\angle ABC\)，故 \(\angle EDC=\angle EBA\).于是 \(\angle ECD=\angle EBA\)，结合 \(E\)，\(C\)，\(B\) 共线可得 \(CD\parallel AB\).
\item 由 \(CD\parallel AB\)，且 \(A\)，\(D\)，\(E\) 共线、\(B\)，\(C\)，\(E\) 共线，三角形 \(EDC\) 与三角形 \(EAB\) 相似.因 \(ED=EC\)，对应边相等，得 \(EA=EB\).

建立直角坐标系，使 \(CD\) 为 \(x\) 轴.由 \(EC=ED\)，可设
\(\displaystyle E=(0,h)\)，\(D=(-c,0)\)，\(C=(c,0)\).由 \(EA=EB\) 且 \(AB\parallel CD\)，可设
\(\displaystyle A=(-u,v)\)，\(B=(u,v)\).又 \(F\)，\(D\)，\(C\)，\(G\) 共线且 \(EF=EG\)，可设
\(\displaystyle F=(-w,0)\)，\(G=(w,0)\).
令 \(\displaystyle \mu=-w^2\)，\(\lambda=\frac{w^2-u^2-v^2}{v}\).则
\(\displaystyle x^2+y^2+\lambda y+\mu=0\) 是一个圆的方程.对 \(F\)，\(G\)，左边分别为 \(w^2+\mu=0\)；对 \(A\)，\(B\)，左边均为 \(u^2+v^2+\lambda v+\mu=0\).因此 \(A\)，\(B\)，\(F\)，\(G\) 都在此圆上，故四点共圆.
\end{enumerate}
\end{solution}

\begin{problem}
在平面直角坐标系 \(xOy\) 中，曲线 \(C_{1}\) 的参数方程为 \(\begin{cases}
x = \cos \varphi,\\
y = \sin \varphi,
\end{cases}\)（\(\varphi\) 为参数）. 曲线 \(C_{2}\) 的参数方程为 \(\begin{cases}
x = a\cos \varphi,\\
y = b\sin \varphi,
\end{cases}\)（\(a > b > 0\)，\(\varphi\) 为参数）. 在以 \(O\) 为极点，\(x\) 轴的正半轴为极轴的极坐标系中，射线 \(l: \theta = \alpha\) 与 \(C_{1}\)，\(C_{2}\) 各有一个交点. 当 \(\alpha = 0\) 时，这两个交点间的距离为 \(2\)，当 \(\alpha = \frac{\pi}{2}\) 时，这两个交点重合.
\begin{enumerate}
\item 分别说明 \(C_{1}\)，\(C_{2}\) 是什么曲线，并求出 \(a\) 与 \(b\) 的值；
\item 设当 \(\alpha = \frac{\pi}{4}\) 时，\(l\) 与 \(C_{1}\)，\(C_{2}\) 的交点分别为 \(A_{1}\)，\(B_{1}\)，当 \(\alpha = -\frac{\pi}{4}\) 时，\(l\) 与 \(C_{1}\)，\(C_{2}\) 的交点分别为 \(A_{2}\)，\(B_{2}\)，求四边形 \(A_{1}A_{2}B_{2}B_{1}\) 的面积.
\end{enumerate}
\end{problem}

\begin{solution}
\begin{enumerate}
\item \(C_1\) 满足 \(x^2+y^2=1\)，所以是以 \(O\) 为圆心、半径 \(1\) 的圆；\(C_2\) 满足
\(\displaystyle \frac{x^2}{a^2}+\frac{y^2}{b^2}=1\)，所以是椭圆.射线 \(\theta=0\) 与两曲线的交点为 \((1,0)\)、\((a,0)\)，距离为 \(2\)，故 \(a=3\).射线 \(\theta=\tfrac{\pi}{2}\) 与两曲线的交点为 \((0,1)\)、\((0,b)\)，两点重合，故 \(b=1\).
\item 此时 \(C_1:x^2+y^2=1\)，\(C_2:\tfrac{x^2}{9}+y^2=1\).当 \(\alpha=\tfrac{\pi}{4}\) 时，令两交点的横坐标分别为 \(r\)，\(s\)，则
\(\displaystyle r=\frac{\sqrt2}{2}\)，\(\frac{s^2}{9}+s^2=1\)，所以 \(\displaystyle s=\frac{3}{\sqrt{10}}\).当 \(\alpha=-\tfrac{\pi}{4}\) 时，交点分别是关于 \(x\) 轴的对称点，四边形 \(A_1A_2B_2B_1\) 是两条竖直边长分别为 \(2r\)，\(2s\)、间距为 \(s-r\) 的梯形.因此面积为
\(\displaystyle \frac{2r+2s}{2}(s-r)=s^2-r^2=\frac9{10}-\frac12=\frac25\).
\end{enumerate}
\end{solution}

\begin{problem}
已知函数 \(f(x) = |x - 2| - |x - 5|\).
\begin{enumerate}
\item 证明： \(-3 \leqslant f(x) \leqslant 3\)；
\item 求不等式 \(f(x) \geqslant x^{2} - 8x + 15\) 的解集.
\end{enumerate}
\end{problem}

\begin{solution}
\begin{enumerate}
\item 由三角不等式
\(\displaystyle |x-2|\leq|x-5|+3\)，得 \(f(x)\leq3\)；由
\(\displaystyle |x-5|\leq|x-2|+3\)，得 \(f(x)\geq-3\).所以 \(-3\leq f(x)\leq3\).
\item 按 \(x=2\)，\(5\) 分段：
\(\displaystyle f(x)=\begin{cases}
-3,&x\leq2,\\
2x-7,&2<x<5,\\
3,&x\geq5.
\end{cases}\)
当 \(x\leq2\) 时，不等式化为 \(x^2-8x+18\leq0\)，判别式为 \(-8<0\)，无解.
当 \(2<x<5\) 时，不等式化为 \(x^2-10x+22\leq0\)，故 \(5-\sqrt3\leq x\leq5+\sqrt3\)，与本区间交得 \([5-\sqrt3,5)\).
当 \(x\geq5\) 时，不等式化为 \(x^2-8x+12\leq0\)，故 \(2\leq x\leq6\)，与本区间交得 \([5,6]\).合并得解集 \([5-\sqrt3,6]\).
\end{enumerate}
\end{solution}
